\documentclass[11pt,letterpaper]{article}
\usepackage[T1]{fontenc}
\usepackage{lmodern}
\usepackage[margin=0.75in,headheight=14pt,headsep=18pt,footskip=24pt]{geometry}
\usepackage{amsmath,amssymb,mathtools}
\usepackage{microtype,xcolor,fancyhdr}
\usepackage[hidelinks]{hyperref}

\definecolor{ink}{HTML}{243E53}
\hypersetup{
  pdftitle={Matrix Elements of a Linear Operator - Note 5},
  pdfauthor={Anonymous},
  pdfsubject={Quantum mechanics notes: orthonormal bases and operator matrices},
  pdfkeywords={quantum mechanics, matrix elements, completeness, linear operators}
}
\newcommand{\ket}[1]{\lvert #1\rangle}
\newcommand{\bra}[1]{\langle #1\rvert}
\newcommand{\braket}[2]{\langle #1\mid #2\rangle}
\newcommand{\mel}[3]{\langle #1\rvert #2\lvert #3\rangle}
\newcommand{\smallheading}[1]{\par\vspace{9pt}{\sffamily\bfseries\color{ink}#1}\par\vspace{3pt}}
\setlength{\parindent}{0pt}
\setlength{\parskip}{5pt}
\setlength{\abovedisplayskip}{7pt}
\setlength{\belowdisplayskip}{7pt}
\pagestyle{fancy}
\fancyhf{}
\fancyhead[L]{\small\sffamily\color{ink}QUANTUM MECHANICS NOTES}
\fancyhead[R]{\small\sffamily\color{ink}NOTE 5}
\fancyfoot[C]{\small 5}
\renewcommand{\headrulewidth}{0.3pt}
\renewcommand{\footrulewidth}{0pt}

\begin{document}
{\Large\sffamily\bfseries\color{ink}1.6\quad Matrix Elements of a Linear Operator}\par
\vspace{5pt}

Let $\{\ket{i}\}_{i=1}^{n}$ be a complete orthonormal basis, and let
$\ket{\psi}=\Omega\ket{v}$. Expanding each state in this basis gives
\[
  \ket{\psi}=\sum_i\ket{i}\braket{i}{\psi},
  \qquad
  \ket{v}=\sum_j\ket{j}\braket{j}{v}.
\]

\smallheading{Expanding the matrix element}
Substituting these expansions into $\mel{w}{\Omega}{v}$, we obtain
\begin{align*}
  \mel{w}{\Omega}{v}
    &= \braket{w}{\psi} \\
    &= \bra{w}\Bigl(\sum_i\ket{i}\braket{i}{\psi}\Bigr) \\
    &= \sum_i\braket{w}{i}\braket{i}{\psi} \\
    &= \sum_i\braket{w}{i}\mel{i}{\Omega}{v} \\
    &= \sum_i\braket{w}{i}\bra{i}\Omega
       \Bigl(\sum_j\ket{j}\braket{j}{v}\Bigr) \\
    &= \sum_{i,j}\braket{w}{i}\mel{i}{\Omega}{j}\braket{j}{v}.
\end{align*}

\smallheading{Matrix representation}
The quantities $\Omega_{ij}=\mel{i}{\Omega}{j}$ are the matrix elements of
$\Omega$ in the chosen basis. The same expression is a matrix product:
\[
\mel{w}{\Omega}{v}=
\begin{bmatrix}
  \braket{w}{1}\\
  \braket{w}{2}\\
  \vdots\\
  \braket{w}{n}
\end{bmatrix}^{\!T}
\begin{bmatrix}
  \mel{1}{\Omega}{1} & \mel{1}{\Omega}{2} & \cdots & \mel{1}{\Omega}{n}\\
  \mel{2}{\Omega}{1} & \mel{2}{\Omega}{2} & \cdots & \mel{2}{\Omega}{n}\\
  \vdots & \vdots & \ddots & \vdots\\
  \mel{n}{\Omega}{1} & \mel{n}{\Omega}{2} & \cdots & \mel{n}{\Omega}{n}
\end{bmatrix}
\begin{bmatrix}
  \braket{1}{v}\\
  \braket{2}{v}\\
  \vdots\\
  \braket{n}{v}
\end{bmatrix}.
\]
Here $T$ denotes the ordinary transpose: the entries $\braket{w}{i}$
already satisfy $\braket{w}{i}=\braket{i}{w}^{*}$.

\smallheading{Completeness: the resolution of the identity}
The basis expansion of an arbitrary state implies
\[
  \boxed{\displaystyle\sum_i\ket{i}\bra{i}=I.}
\]
Thus the double sum can also be obtained by inserting the identity on
both sides of the operator:
\begin{align*}
  \mel{w}{\Omega}{v}
    &= \bra{w}\Bigl(\sum_i\ket{i}\bra{i}\Bigr)\Omega
       \Bigl(\sum_j\ket{j}\bra{j}\Bigr)\ket{v}\\
    &= \sum_{i,j}\braket{w}{i}\mel{i}{\Omega}{j}\braket{j}{v}.
\end{align*}
\end{document}
